\documentclass[11pt]{article}
\usepackage[a4paper,margin=28mm]{geometry}
\usepackage[T1]{fontenc}
\usepackage[utf8]{inputenc}
\usepackage{lmodern}
\usepackage{amsmath,amssymb,amsthm,mathtools}
\usepackage{booktabs}
\usepackage{array}
\usepackage{xcolor}
\usepackage{hyperref}
\hypersetup{colorlinks=true,linkcolor=blue!45!black,citecolor=blue!45!black,urlcolor=blue!55!black}
\usepackage{microtype}

\newtheorem{theorem}{Theorem}[section]
\newtheorem{proposition}[theorem]{Proposition}
\newtheorem{lemma}[theorem]{Lemma}
\newtheorem{remark}[theorem]{Remark}
\newcommand{\Fp}{\mathbf F_p}
\newcommand{\Status}[1]{\par\smallskip\noindent\textbf{Status: #1.}\ }

\title{\textbf{Arithmetic Quantum Cat Maps:}\\
Exact Moment Reductions, Finite Certificates, and the Remaining BGS Obstruction}
\author{Maksim Churilov}
\date{Research draft --- 3 September 2026}

\begin{document}
\maketitle

\begin{abstract}
We give a conservative reconstruction of a long computational programme on
arithmetic quantum cat maps. The main verified contribution is an exact
pair-key reduction for the fourth-moment geometry. It isolates a weighted
diagonal key, decomposes the resulting difference distribution into off/off,
mixed, and diagonal pieces, and converts the mixed term through a Cayley
transform into an additive-energy problem. An exhaustive integer matching
audit and exact cyclotomic resultants give the exceptional prime sets
\[
E_{11}=\{23,67,89\},\qquad
E_{13}=\{53,79,131,157,313\},
\]
while the exact design $(p,K)=(461,23)$ has collision energy $1190$ against a
forced baseline $902$, ruling out a naive uniform-in-$K$ extrapolation.
The remaining off/off term admits a promising elliptic-family organization,
but the required sheaf, isogeny, bad-fiber, and boundary hypotheses are not
yet closed; the associated square-root correlation estimate is therefore a
theorem draft, not a theorem. We also record the scope of earlier finite Lean
certificates, twisted-Jacobi regressions, antiunitary symmetry calculations,
and finite GOE comparisons. In particular, no result in this paper proves
microscopic Bohigas--Giannoni--Schmit (BGS) eigenphase statistics.
\end{abstract}

\section{Purpose, conventions, and evidence levels}

The project began as an attempt to connect arithmetic structure in quantized
hyperbolic toral automorphisms with random-matrix spectral statistics. The
historical workspace accumulated correct finite identities, useful numerical
experiments, incomplete proof routes, and several false ``proved'' labels.
This manuscript states only the part that survives a reconstruction audit.

Throughout, $p$ denotes an odd prime and $H\leq\Fp^\times$ a cyclic subgroup
of order $K$. Division denotes multiplication by the inverse in $\Fp$. When a
statement concerns the audited designs $K=11,13,23$, this is stated
explicitly. We use four evidence levels:
\begin{itemize}
\item \textbf{Exact}: a current algebraic derivation.
\item \textbf{Finite certificate}: an exhaustive or machine-checked statement
      at explicitly specified finite parameters.
\item \textbf{Numerical}: floating-point evidence only.
\item \textbf{Draft}: a proof architecture with named unresolved obligations.
\end{itemize}
A finite certificate is not silently upgraded to a uniform theorem, and a
numerical fit is not a hypothesis test unless its statistical assumptions are
verified.

The relevant established cat-map literature is mainly about eigenfunctions,
Hecke symmetries, quantum ergodicity, and matrix-element fluctuations.
Kurlberg--Rudnick compute second and fourth moments for normalized matrix
elements and compare them with a conjectural SU(2)-trace law \cite{KR2005};
short-window results concern sums of matrix coefficients \cite{KRR2007}.
Recent higher-dimensional quantum-ergodicity results \cite{KORS2025,BS2026}
do not provide a microscopic eigenphase local law. This distinction is
load-bearing for the interpretation below.

\section{The exact pair-key geometry}

Let $x,y\in\Fp\setminus\{0,1\}$ and set
\begin{equation}\label{eq:keydef}
 u=\frac{x}{y},\qquad
 v=\frac{x-1}{y-1},\qquad
 d=x-y.
\end{equation}

\begin{proposition}[Off-diagonal pair-key bijection]\label{prop:key}
If $x\neq y$, then $u,v\in\Fp^\times\setminus\{1\}$, $u\neq v$, and
\begin{equation}\label{eq:inverse}
 y=\frac{1-v}{u-v},\qquad
 x=\frac{u(1-v)}{u-v},\qquad
 d=\frac{(u-1)(v-1)}{v-u}.
\end{equation}
Conversely every pair $u,v\in\Fp^\times\setminus\{1\}$ with $u\neq v$
arises from exactly one off-diagonal pair $(x,y)$.
\end{proposition}

\begin{proof}
From $x=uy$ and $x-1=v(y-1)$ we obtain $(u-v)y=1-v$, which gives the
first two formulas in \eqref{eq:inverse}. Then
\[
x-y=(u-1)y
     =\frac{(u-1)(1-v)}{u-v}
     =\frac{(u-1)(v-1)}{v-u}.
\]
The domain restrictions are immediate from \eqref{eq:keydef}; the displayed
inverse proves injectivity and surjectivity.
\end{proof}

The diagonal is exceptional in a useful and completely explicit way: all
$p-2$ pairs $x=y\in\Fp\setminus\{0,1\}$ collapse to the single key
\[
(u,v,d)=(1,1,0)
\]
with multiplicity $p-2$. Any pair-key calculation that treats this key with
multiplicity one is wrong. This weighted diagonal is the mechanism behind
several historical factor errors.

\section{Fourth-moment coupling and the exact distribution}

Let a collapsed key carry its multiplicity $n(u,v)$, equal to one off the
diagonal and $p-2$ at $(1,1)$. Two ordered keys $(u_1,v_1)$ and $(u_2,v_2)$
are coupled when
\[
u_1u_2\in H,\qquad v_1v_2\in H.
\]
Let $B(t)$ be the multiplicity-weighted number of coupled ordered key pairs
whose differences sum to $t$.

Define $O(t)$ to be the off/off contribution and
\[
C_H(t)=\#\left\{(u,v)\in(H\setminus\{1\})^2:
u\neq v,\ \frac{(u-1)(v-1)}{v-u}=t\right\}.
\]
Then the diagonal bookkeeping is exact:
\begin{equation}\label{eq:Bdecomp}
 B(t)=O(t)+2(p-2)C_H(t)+(p-2)^2\mathbf 1_{t=0}.
\end{equation}
The factor $2$ is forced because diagonal/off and off/diagonal are distinct
ordered classes.

\begin{proposition}[Exact total coupled count]\label{prop:v4}
For the corrected pair-key model,
\begin{equation}\label{eq:v4}
 |V_4|=(K^2+1)p^2-(6K^2+3K)p+13K^2-3K.
\end{equation}
\end{proposition}

\Status{Exact formula with finite cross-checks.}
The repository verifies \eqref{eq:v4} by independent implementations. At
$(p,K)=(23,11)$ the total is $48621$, while at $(53,13)$ it is $423879$.
The Lean certificate at $(53,13)$ independently checks the diagonal
multiplicity, off-diagonal key injectivity, and the latter $V_4$ count.

\section{Character orthogonality and Fourier moments}

For the audited character convention, orthogonality collapses the
fourth-moment character sum to the finite Fourier transform of $B$:
\begin{equation}\label{eq:fourier}
 \sum_{\chi,\psi}|J(\chi,\psi;q)|^4
 =m^2\sum_{t\in\Fp}B(t)e_p(qt),
\end{equation}
where $e_p(z)=\exp(2\pi iz/p)$ and $m$ is the annihilator dimension in the
specified design.

Equation \eqref{eq:fourier} is important because the averaging in $q$ is then
deterministic finite Fourier algebra, not a probabilistic independence
assumption. If
\[
S(q)=\sum_t B(t)e_p(qt),
\]
then the mean over $q\neq0$ follows from
\[
\sum_{q\neq0}e_p(qt)=
\begin{cases}
p-1,&t=0,\\
-1,&t\neq0,
\end{cases}
\]
and Parseval determines the centered second moment from the $L^2$ energy of
the nonzero-$t$ part of $B$. Thus the fluctuation problem is reduced to
controlling the combinatorial energies in \eqref{eq:Bdecomp}.

\section{The mixed term as Cayley additive energy}

Introduce the Cayley transform
\[
c(h)=\frac{h+1}{h-1},\qquad h\in H\setminus\{1\}.
\]
A direct calculation gives
\[
c(u)-c(v)=\frac{2(v-u)}{(u-1)(v-1)},
\]
hence
\begin{equation}\label{eq:cayley}
d(u,v)=\frac{2}{c(u)-c(v)}.
\end{equation}
Let $A=c(H\setminus\{1\})$. Since $x\mapsto2/x$ permutes $\Fp^\times$,
\eqref{eq:cayley} identifies the nonzero multiplicities of $C_H$ with the
nonzero difference multiplicities of $A$. Therefore
\begin{equation}\label{eq:energy}
 \sum_t C_H(t)^2=E_+(A)-|A|^2,
\end{equation}
where $E_+(A)=\sum_x r_{A-A}(x)^2$.

Moreover $c(h^{-1})=-c(h)$, so $A=-A$. This central symmetry forces a
baseline of repeated differences. With $|A|=K-1$ and $0\notin A$, the
collision baseline used by the finite verifier is
\begin{equation}\label{eq:baseline}
 (K-1)(2K-5).
\end{equation}

The same differences can be written as intersections of $H\setminus\{1\}$
with a fractional-linear image, giving an exact PGL$_2$ interpretation of the
mixed term. This is useful structurally, but the old attempt to apply a
uniform square-root bound directly to the historical M\"obius fibers was
false: the $d=0$ fiber degenerates to a constant map.

\section{Exact matching audit and cyclotomic exceptional primes}

The collision equation between two Cayley differences can be cleared of
denominators and expressed in powers of a primitive $K$-th root. Matching
the two six-term multisets gives $6!=720$ integer linear systems in the
exponent variables.

\begin{proposition}[Finite integer matching classification]\label{prop:matching}
The active exact verifier classifies the 720 systems as
\[
30\text{ of rank }2,\qquad
240\text{ of rank }3,\qquad
450\text{ of rank }4.
\]
The gcds of maximal nonzero minors show that the only rational prime at which
these integer matrices can acquire an additional rank drop is $2$.
Among the rank-$2$ and rank-$3$ kernels, the admissible positive-dimensional
systems are exhausted by the identity and inversion-swap symmetries; the
others force a forbidden zero or diagonal configuration.
\end{proposition}

\Status{Exact exhaustive finite algebra.}
This is recomputed in CI from the integer matrices rather than accepted from
a stored JSON file.

For a specified odd prime $K$, a normalized non-forced collision gives a
cyclotomic polynomial value in $\mathbb Z[\zeta_K]$. When this algebraic
integer is nonzero, a finite-field collision can occur only at rational primes
dividing its norm. The current repository performs the complete normalized
nonvanishing/resultant enumeration for $K=11$ and $K=13$.

\begin{theorem}[Audited exceptional sets]\label{thm:exceptions}
The exact resultant computation gives
\[
E_{11}=\{23,67,89\},\qquad
E_{13}=\{53,79,131,157,313\}.
\]
For $K=11$ the enumeration checks $793$ normalized non-forced quadruples and
finds $8$ distinct nonzero norms; for $K=13$ it checks $1431$ quadruples and
finds $14$ distinct nonzero norms.
\end{theorem}

The finite-field regression tests both directions of the prediction:
\[
\begin{array}{c|c|c|c}
(p,K)&\text{exceptional?}&
\sum_t C_H(t)^2-\text{baseline}&\text{result}\\\hline
(53,13)&\text{yes}&208&\text{excess observed}\\
(67,11)&\text{yes}&48&\text{excess observed}\\
(199,11)&\text{no}&0&\text{forced only}\\
(521,13)&\text{no}&0&\text{forced only}.
\end{array}
\]
The historical verifier had a weaker one-way condition in which an expected
exceptional prime passed without requiring observed excess; the active
verifier closes that loophole.

\subsection{A non-uniformity witness}

At $(p,K)=(461,23)$ the exact finite collision energy is
\[
\sum_t C_H(t)^2=1190,
\]
whereas \eqref{eq:baseline} equals $22(46-5)=902$. The excess is $288$.
This does not contradict the fixed-parameter cyclotomic mechanism; it shows
that a statement calibrated at $K=11,13$ cannot be silently promoted to a
uniform-in-$K$ bound.

\section{The off/off elliptic route}

After the exact mixed term is separated, the unresolved part of the moment
problem is the centered energy of $O(t)$ in \eqref{eq:Bdecomp}. Historical
symbolic computations organize the relevant fibers as a one-parameter
genus-one family and produce a nonconstant rational $j$-map, visible
degenerate fibers, and same-$j$ parameter symmetries.

\Status{Theorem draft, not a theorem.}
Nonconstant $j$ proves only non-isotriviality of the candidate generic family.
It does \emph{not} by itself imply the trace-function correlation bound needed
here. A publication-grade argument must still establish:
\begin{enumerate}
\item a smooth projective model over an explicitly specified open base;
\item the complete discriminant and bad-fiber locus, including characteristic
      exceptions;
\item the relation between the affine counting problem and the deleted
      projective boundary;
\item purity, lissity, geometric irreducibility, and a uniform
      conductor/ramification bound for the actual $R^1$ sheaf;
\item all geometric isomorphisms and twists between parameter fibers;
\item the relevant separable and inseparable isogeny exclusions;
\item the exact hypotheses of the chosen sheaf-correlation theorem.
\end{enumerate}

Michel-type family arguments \cite{Michel1995}, quantitative sheaf theory
\cite{SFFK2023}, Tate/isogeny technology such as \cite{Zarhin2013}, and the
recent joint-family example of Grove--Saad \cite{GroveSaad2025} are relevant
tools and templates. They are not plug-in proofs for this pencil.

If the required square-root correlation step is established with uniform
constants, the existing combinatorics lead to the target
\begin{equation}\label{eq:target}
 O\!\left(K^2p+K^4p^{1/2}+K^4\right).
\end{equation}
The error in \eqref{eq:target} is lower order on the intended main scale in
the regime $K=o(p^{1/4})$. Both \eqref{eq:target} and this regime remain
conditional until the listed geometric obligations are closed.

\section{Earlier results retained after the audit}

A central aim of the reconstruction was to avoid the opposite errors of
trusting every historical ``PROVEN'' label and deleting every result not on
the newest proof route.

\subsection{Finite torus geometry at p=199}

The active Lean module \texttt{TorusLevelSets.lean} verifies for
\[
Q(a,b)=a^2-ab-b^2\pmod{199}
\]
and the cyclic action of
$A^T=\left(\begin{smallmatrix}2&1\\1&1\end{smallmatrix}\right)$:
the action generator has order $11$; every nonzero level has $198$ nonzero
points; the zero cone minus the origin has $396$ points; every nonzero level
has $18$ cyclic orbits; and the action is free on those nonzero levels. These
are finite \texttt{native\_decide} certificates. The analytic Weil
intertwining statement that motivated them is not formalized by this file.

\subsection{Finite branch scale at p=199}

The active Lean module \texttt{TorusScaleLaw.lean} checks that $a=139$ has
order $11$, the recorded inverse is $63$, the half-index trace is $123$, and
the branch scales $123$ and $76=-123$ satisfy
\[
123^2\equiv76^2\equiv5,\qquad 123\cdot76\equiv-5\pmod{199}.
\]
Again, this is a finite design certificate, not a general scale theorem.

\subsection{Twisted-Jacobi finite regressions}

The preserved scalar/vectorized checks and $q$-curve moment records verify
the implementation at the stated designs. A historical domain omission was
identified: the old closed form requires the relevant $n_1\neq0$; on the
audited nonzero-$q$ curves this is automatic, but the unrestricted old
statement is not retained.

\subsection{Antiunitary symmetry: the safe statement}

If an antiunitary $\Theta$ obeys $\Theta U\Theta^{-1}=U^{-1}$ and
$Uv=\lambda v$ with $|\lambda|=1$, then anti-linearity gives
\[
U\Theta v=\bar\lambda^{-1}\Theta v=\lambda\Theta v.
\]
Thus the relevant eigenspace is preserved. An earlier objection that
$\Theta$ necessarily sends it to a distinct $\lambda^{-1}$ sector forgot the
conjugation from anti-linearity.

What is \emph{not} promoted is a blanket claim that every relevant global
block has already been certified as GOE class. The historical finite/kernel
chart remains useful; a full blockwise statement requires a fresh treatment
of all charts and conventions if it becomes load-bearing.

\subsection{Shape, Catalan, and arcsine material}

The archive contains exact finite relation counts and moment patterns that
motivated Catalan/arcsine conjectures. They remain useful for discovery and
regression. We do not promote the historical all-orders full-group arcsine
package because its word-trace reduction has an unresolved higher rewrite
case. Agreement of finitely many moments is not an all-orders proof.

\section{Finite spectral diagnostics}

Finite parity-reduced spectra at the recorded designs show descriptive
agreement with random-matrix benchmarks. The current repository deliberately
downgrades two historical statistical labels.

First, the spacing record now reports only the empirical CDF sup-distance to
the GOE Wigner surmise. The historical script compared this number with an
IID Kolmogorov--Smirnov threshold, but gaps obtained after polynomial
unfolding and pooled within the same spectra are dependent. We therefore do
not report a KS ``pass''.

Second, the old file called ``number variance'' used forward windows starting
at each eigenvalue and truncated at the spectral edge. This is not the
standard stationary number variance $\Sigma^2(L)$. The values are retained as
a matched-GOE window-count diagnostic because the cat-map and GOE samples
were processed by the same algorithm, but the old two-sigma pass/fail
interpretation is withdrawn.

These numerical observations are compatible with the random-matrix
motivation and with earlier numerical cat-map studies \cite{GLR2003}; they
are not evidence for a limiting local law.

\section{Why the BGS problem remains open}

The BGS claim concerns microscopic eigenphase correlations after all relevant
symmetries are removed. The exact pair-key reduction and even the prospective
estimate \eqref{eq:target} are moment statements about a structured arithmetic
family. A local-statistics theorem needs a separate mechanism that controls
eigenvalue correlations at the mean-spacing scale. No such transfer theorem
has been verified for the present family.

Accordingly, the current logical chain is
\[
\begin{array}{c}
\boxed{\text{exact pair-key geometry}}
\longrightarrow
\boxed{\text{mixed Cayley energy}}
\\[4pt]
\Downarrow\\[-2pt]
\boxed{\text{off/off elliptic correlation: open}}
\longrightarrow
\boxed{\text{moment theorem}}
\\[4pt]
\Downarrow\\[-2pt]
\boxed{\text{microscopic BGS bridge: open}}
\end{array}
\]
Closing an earlier box does not silently close a later one.

\section{Reproducibility and negative controls}

The active repository keeps a small scientific surface:
\begin{itemize}
\item Python exact regressions and finite-field primitives;
\item two active C++ programs: an independent fourth-moment pair-key verifier
      and an exact cat-map order/split-prime utility;
\item six active Lean finite-certificate modules;
\item promoted finite certificates and explicitly descriptive numerical data;
\item a separate historical record for retractions and old experiments.
\end{itemize}

The $K=11,13$ Cayley certificate is regenerated in CI and compared
structurally with the committed JSON. The C++ verifier independently checks
the pair-key geometry at finite designs. CI also checks that every active Lean
module is imported before running \texttt{lake build}.

Negative controls are retained when they explain a failure mode. In
particular, the degenerate historical M\"obius fiber is formalized as a finite
falsifier, and $(461,23)$ is retained as a non-uniformity witness. Historical
Julia source is inventoried for later migration; false algorithms are not
translated merely to preserve a language port.

\section{Retractions that affect interpretation}

The following historical conclusions are explicitly not part of the present
result:
\begin{itemize}
\item the M\"obius/M4D proof of the mixed/off-diagonal fluctuation bound;
\item the one-sided mixed-class count that omitted one of the two ordered
      diagonal/off classes;
\item any shortcut from a global random-matrix universality theorem to this
      structured deterministic arithmetic family without verified hypotheses;
\item the historical unrestricted twisted-Jacobi formula that omitted its
      domain condition;
\item the all-orders arcsine/Catalan conclusion based on the unresolved
      word-trace rewrite;
\item finite numerical GOE ``passes'' interpreted as proof of a limiting law.
\end{itemize}
These failures are kept in the research record because they are useful
regression tests; they are excluded from the active theorem surface.

\section{Conclusion}

The reconstruction leaves a substantially narrower but stronger result than
the historical project claimed. The pair-key geometry, weighted diagonal,
fourth-moment coupling, Fourier reduction, Cayley transform, integer matching
audit, the explicit $E_{11}$ and $E_{13}$ resultant sets, and the $(461,23)$
obstruction are reproducible. Earlier finite torus, scale, and twisted-Jacobi
results survive with their actual finite scope.

The main analytic task is now sharply localized: turn the off/off elliptic
organization into a theorem by verifying the precise sheaf and isogeny
hypotheses. Even after that, microscopic BGS statistics require a separate
local-law bridge. Keeping these closure conditions explicit is more useful
than a stronger but unsupported theorem statement.

\begin{thebibliography}{99}

\bibitem{KR2001}
P.~Kurlberg and Z.~Rudnick,
\emph{Value distribution for eigenfunctions of desymmetrized quantum maps},
arXiv:math-ph/0101008 (2001).

\bibitem{KR2005}
P.~Kurlberg and Z.~Rudnick,
\emph{On the distribution of matrix elements for the quantum cat map},
Ann. of Math. \textbf{161} (2005), 489--507,
doi:10.4007/annals.2005.161.489.

\bibitem{KRR2007}
P.~Kurlberg, L.~Rosenzweig and Z.~Rudnick,
\emph{Matrix elements for the quantum cat map: fluctuations in short windows},
arXiv:math/0701685 (2007).

\bibitem{GLR2003}
A.~Gamburd, J.~Lafferty and D.~Rockmore,
\emph{Eigenvalue spacings for quantized cat maps},
J. Phys. A \textbf{36} (2003), 3487--3499,
doi:10.1088/0305-4470/36/12/336.

\bibitem{KORS2025}
P.~Kurlberg, A.~Ostafe, Z.~Rudnick and I.~E.~Shparlinski,
\emph{On Quantum Ergodicity for Higher Dimensional Cat Maps},
Commun. Math. Phys. \textbf{406} (2025), 174,
doi:10.1007/s00220-025-05350-1.

\bibitem{BS2026}
S.~Bhakta and I.~E.~Shparlinski,
\emph{On quantum ergodicity for higher-dimensional cat maps modulo prime powers},
Bull. London Math. Soc. \textbf{58} (2026),
doi:10.1112/blms.70300.

\bibitem{Michel1995}
P.~Michel,
\emph{Rang moyen de familles de courbes elliptiques et lois de Sato--Tate},
Monatsh. Math. \textbf{120} (1995), 127--136.

\bibitem{SFFK2023}
W.~Sawin, A.~Forey, J.~Fres\'an and E.~Kowalski,
\emph{Quantitative sheaf theory},
J. Amer. Math. Soc. \textbf{36} (2023), 653--726,
doi:10.1090/jams/1008.

\bibitem{Zarhin2013}
Yu.~G.~Zarhin,
\emph{Homomorphisms of abelian varieties over geometric fields of finite characteristic},
J. Inst. Math. Jussieu \textbf{12} (2013), 225--236,
doi:10.1017/S1474748012000655.

\bibitem{GroveSaad2025}
B.~Grove and H.~Saad,
\emph{Hypergeometric distributions and joint families of elliptic curves},
Res. Math. Sci. \textbf{12} (2025), 74,
doi:10.1007/s40687-025-00554-2.

\end{thebibliography}

\end{document}
